\chapter{Start here}
\label{ch:start}

\section{Who this guide is for}

You want to use BigMath to work with integers that are too big for a machine word:
cryptographic sizes, exact combinatorics, computing millions of digits of a constant,
or a calculator that never overflows. You do not need to know how big-number
arithmetic works inside. But knowing it helps you predict how long an operation will
take, pick the right API, and understand what the build options do.

This guide explains every algorithm the library uses, in plain language. Each one is
explained in more than one way: a picture, a small example you can check by hand,
and the reason it is faster than the simpler method. It also explains \emph{when} the
library switches from one algorithm to another, and why the switch points sit where
they do.

You will not need more mathematics than school algebra. Where a formula appears, the
text around it says in words what it means.

\section{How to read it}

The chapters build on each other, but you can jump around:

\begin{itemize}
  \item \Cref{ch:repr,ch:basics} cover how numbers are stored and how the small
        building blocks (add, subtract, multiply by one limb) work. Everything
        else is made of these.
  \item \Cref{ch:mul,ch:ntt,ch:fermat} explain multiplication from schoolbook
        through the NTT and the experimental Fermat FFT. See \cref{ch:muldispatch}
        for the dispatch rules.
  \item \Cref{ch:div,ch:fastdiv,ch:divdispatch} cover division, including
        Hensel exact division (\cref{sec:exactdivision}).
  \item \Cref{ch:conv} covers decimal strings, radix conversion (\cref{sec:convertbase}),
        and fixed-point arithmetic.
  \item \Cref{ch:numbertheory} covers GCD, integer roots, modular powers, prime
        searches, combinatorics, and reconstruction from modular residues.
  \item \Cref{ch:parallel} covers the thread pool.
  \item \Cref{ch:practice} is a checklist for using the library well, and
        \cref{ch:decisions} records the ideas that were tried and rejected.
  \item \Cref{app:ref} lists every threshold in one place.
\end{itemize}

Six kinds of boxes appear throughout:

\begin{keyidea}
The one sentence to remember about a section.
\end{keyidea}
\begin{example}
A small calculation you can follow with pencil and paper.
\end{example}
\begin{another}
The same idea explained from a different angle.
\end{another}
\begin{perfnote}
What makes the method fast on a real processor.
\end{perfnote}
\begin{practice}
Advice for someone calling the library.
\end{practice}
\begin{pitfall}
Something that went wrong during development, and what it taught.
\end{pitfall}

\section{Words used in this guide}

\begin{description}[leftmargin=2.2cm,style=nextline]
  \item[limb] One machine word of a big number. BigMath limbs are 64 bits wide,
        so a limb holds a ``digit'' between $0$ and $2^{64}-1$. A 1000-digit decimal
        number needs about 52 limbs.
  \item[$B$] The base of the limbs, $B = 2^{64}$. A number with limbs
        $a_0, a_1, \dots, a_{n-1}$ has the value
        $a_0 + a_1 B + a_2 B^2 + \dots + a_{n-1}B^{n-1}$.
  \item[$n$, $n_a$, $n_b$] Lengths in limbs.
  \item[$\M(n)$] The time to multiply two $n$-limb numbers. Division and
        conversion are often described as ``a few $\M(n)$''.
  \item[threshold] A size at which the library switches algorithm. All
        thresholds were measured, not guessed (\cref{sec:measured}).
  \item[dispatcher] The function that looks at the operand sizes and picks the
        algorithm.
  \item[warm / cold pool] Whether the worker threads are awake (warm) or asleep
        (cold) when an operation starts. See \cref{ch:parallel}.
\end{description}

As a rule of thumb, one limb is about 19.3 decimal digits:

\begin{center}
\begin{tabular}{rrr}
\toprule
decimal digits & bits & limbs \\
\midrule
100 & 333 & 6 \\
1\,000 & 3\,322 & 52 \\
10\,000 & 33\,220 & 520 \\
1\,000\,000 & 3.3 million & 52\,000 \\
100\,000\,000 & 332 million & 5.2 million \\
\bottomrule
\end{tabular}
\end{center}

\section{The whole library on one page}
\label{sec:onepage}

\Cref{tab:onepage} is the short version of this guide. For each operation it shows
which algorithm runs at which size with the generic thresholds. Sizes are in
limbs unless marked as bits or decimal digits. Platform profiles can override the thresholds
(\cref{sec:platformprofiles}).

\begin{table}[!t]
\centering
\small
\renewcommand{\arraystretch}{1.10}
\begin{tabularx}{\textwidth}{@{}l l >{\raggedright\arraybackslash}X@{}}
\toprule
operation & size (limbs) & algorithm \\
\midrule
$a \pm b$ & any & serial carry or borrow chain (\cref{sec:paralleladdsub}) \\
\midrule
\multirow{5}{*}{$a \times b$}
  & both $\le 16$ & stack-buffer schoolbook (\cref{sec:smallops}) \\
  & shorter $\le 24$, or total $\le 56$ & schoolbook (\cref{sec:schoolbook}) \\
  & up to $\approx 310$ each & Karatsuba; levels with $\ge 168$ limbs split three ways (Toom-3) (\cref{sec:karatsuba,sec:toom}) \\
  & shorter $\ge 64$, area $\ge 96\,000$ & NTT if the thread pool is awake, else Karatsuba (\cref{sec:poolrule}) \\
  & area $\ge 262\,144$ ($\approx 512\times512$) & NTT always (\cref{ch:ntt}) \\
\midrule
\multirow{3}{*}{$a^2$}
  & $< 352$ & Karatsuba square (Toom-3 from 160) \\
  & $352$--$607$ & NTT if pool awake, else Karatsuba square \\
  & $\ge 608$ & NTT (only three forward transforms) \\
\midrule
\multirow{5}{*}{$a \div b$}
  & divisor 1 limb & short division with a precomputed reciprocal (\cref{sec:shortdiv}) \\
  & divisor $< 96$ & schoolbook long division, Knuth's Algorithm D (\cref{sec:knuth}) \\
  & quotient $\ge$ divisor & Burnikel--Ziegler, or Newton from divisor 1536 (768 if ratio $\ge 3$) \\
  & quotient $<$ divisor & Burnikel--Ziegler; Newton from 3072 at ratio $\ge 7/4$; quotient-sized if the quotient has $\ge 3072$ limbs \\
  & \code{a / b} with short quotient & divide only the top limbs (\cref{sec:toplimbs}) \\
\midrule
parse & $\le 2048$ digits & 18 digits at a time, multiply-and-add (\cref{sec:linparse}) \\
      & larger & split in half, recurse, combine with one multiply; 8 threads from 100\,000 digits \\
\midrule
format & standalone 48--96 & scaled-remainder chunk extraction (\cref{sec:scaledformat}) \\
       & other inputs $<1024$ digits & direct for one limb; 19-digit short division otherwise \\
       & larger inputs & split by powers of ten; scaled or scalar leaves \\
\midrule
GCD & shorter $\ge 65\,536$ bits & half-GCD for comparable widths; extended cutoff 16\,384 bits (\cref{sec:halfgcd}) \\
modular power & odd modulus $\ge 64$ & FastREDC with exponents from 64 bits (\cref{sec:fastredc}) \\
\bottomrule
\end{tabularx}
\caption{Arithmetic dispatch with generic thresholds and 64-bit limbs.}
\label{tab:onepage}
\end{table}

\section{Why the switch points are measured}
\label{sec:measured}

Textbooks say ``Karatsuba is faster than schoolbook for large $n$''. That is true but
not useful: \emph{how} large? The answer depends on the processor, the compiler, and
how carefully each algorithm is written. A clever algorithm has extra bookkeeping. It
only wins once the numbers are big enough for the savings to pay for that bookkeeping.

BigMath measures crossovers on a grid of sizes and shapes. The generic
thresholds come from \file{include/biginteger/build/DispatchThresholds.h}.
A platform profile supplies overrides for its compiler and target. You
can also override individual thresholds at build time, for example with
\code{-DBIGMATH\_KARATSUBA\_THRESHOLD=32}. The timing examples in this
guide are measurements from an Apple M1 Max. Crossovers depend on the
target CPU, compiler, and worker state.

\begin{keyidea}
An algorithm with a better growth rate always wins eventually, but ``eventually'' is
a number that has to be measured. BigMath uses the simplest method that is fastest at
each size, and switches at measured crossover points.
\end{keyidea}
